\subsection{Multiplicities and Lengths in Semisimple Modules}

PROPOSITION 10. --- Let M be a semisimple A-module. Let \((\mathrm{M}_{i})_{i\in\mathrm{I}}\) be a family of simple submodules with direct sum M. The following properties are equivalent:

\begin{enumerate}
\def\labelenumi{(\roman{enumi})}
\item
  M has finite length.
\item
  M is Artinian.
\item
  M is Noetherian.
\item
  M is finitely generated.
\item
  I is finite.
\end{enumerate}

If M has these properties, then the length of M is equal to the cardinal of I.

If the set I is finite, then M has properties (i), (ii), (iii), and (iv). Suppose that the set I is infinite. By Example 2 of VIII, p.~2, the module M is neither Artinian nor Noetherian; because every module of finite length is Artinian and Noetherian (VIII, p.~2, Proposition 1), M also does not have finite length. Finally, every element of M belongs to the sum of a finite number of submodules \(M_{i}\) , so M is not finitely generated. This proves the equivalence of properties (i) through (v). If these hold, then we have \(\text{long}(M) = \sum_{i \in I} \text{long}(M_{i}) = \text{Card}(I)\) (II, §1, No.~10, p.~213, Corollary 5).

PROPOSITION 11. --- Let M be a semisimple A-module that is the direct sum of a family \((\mathrm{M}_{i})_{i\in\mathrm{I}}\) of simple submodules. For every \(\lambda\in S\) , we denote by \(\mathrm{I}(\lambda)\) the set of indices \(i\in I\) such that \(M_{i}\) is of class \(\lambda\) . The cardinal of \(\mathrm{I}(\lambda)\) is equal to the dimension of the left \(D_{\lambda}\) -vector space \(\mathrm{Hom}_{\mathrm{A}}(\mathrm{S}_{\lambda},\mathrm{M})\) .

The isotypical component of A of type \(\lambda\) is isomorphic to \(\mathrm{S}_{\lambda}^{(\mathrm{I}(\lambda))}\) (VIII, p.~65, Proposition 4, b)). The \(D_{\lambda}\) -vector space \(\mathrm{Hom}_{\mathrm{A}}(\mathrm{S}_{\lambda},\mathrm{M})\) may be identified with \(\mathrm{Hom}_{\mathrm{A}}(\mathrm{S}_{\lambda},\mathrm{M}_{\lambda})\) , so is isomorphic to \(\mathrm{D}_{\lambda}^{(\mathrm{I}(\lambda))}\) (No.~2). This proves the proposition.

Every simple module is primordial (VIII, p.~45), so every semisimple module is semiprimordial. Let M be a semisimple A-module, and let \(\lambda \in \mathcal{S}\) . The multiplicity of \(\lambda\) in M is the primordial multiplicity \([\mathrm{M}:\lambda]\) of \(\lambda\) in M defined in VIII, p.~34. Proposition 11 translates to the equality

\[
[ \mathrm{M}: \lambda ] = \dim_ {\mathrm{D} _ {\lambda}} (\operatorname{Hom} _ {\mathrm{A}} (\mathrm{S} _ {\lambda}, \mathrm{M})).\tag{11}
\]

More generally, if \(((\mathrm{V}_{\lambda})_{\lambda \in \mathcal{S}},\alpha)\) is a description of M, then \([\mathrm{M}:\lambda ]\) is equal to \(\dim_{\mathrm{D}_{\lambda}}(\mathrm{V}_{\lambda})\) . By Proposition 6 of VIII, p.~68, we also have

\[
[ \mathrm{M}: \lambda ] = \dim_ {\mathrm{D} _ {\lambda}} (\operatorname{Hom} _ {\mathrm{A}} (\mathrm{M}, \mathrm{S} _ {\lambda}))\tag{12}
\]

when the multiplicity \([M : \lambda]\) is finite. Semisimple A-modules M and \(M'\) are isomorphic if and only if we have \([M : \lambda] = [M' : \lambda]\) for every \(\lambda \in S\) .

Let M be a semisimple A-module. There exists a cardinal I that has the following property: for every decomposition \(M = \bigoplus_{i \in I} M_i\) of M into a direct sum of simple modules, the cardinal of I is equal to I (VIII, p.~34, Corollary 2). This cardinal is called the length of the semisimple A-module M and denoted by \(\operatorname{long}_{\mathrm{A}}(\mathrm{M})\) or \(\operatorname{long}(\mathrm{M})\) . When M has finite length, this definition is compatible with that of II, §1, No.~10, p.~212, by Proposition 10.

The simple A-modules are the semisimple A-modules of length 1, and we have the formula

\[
\operatorname{long} _ {\mathrm{A}} \left(\bigoplus_ {j \in \mathrm{J}} \mathrm{M} _ {j}\right) = \sum_ {j \in \mathrm{J}} \operatorname{long} _ {\mathrm{A}} (\mathrm{M} _ {j})\tag{13}
\]

for every family \((\mathrm{M}_j)_{j\in \mathrm{J}}\) of semisimple A-modules. By Proposition 11, we have

\[
\mathrm{long} _ {\mathrm{A}} (\mathrm{M}) = \sum_ {\lambda \in \mathcal {S}} \dim_ {\mathrm{D} _ {\lambda}} \mathrm{Hom} _ {\mathrm{A}} (\mathrm{S} _ {\lambda}, \mathrm{M}).\tag{14}
\]

By applying this formula to \(M_{\lambda}\) , we obtain \([M : \lambda] = \text{long}_{A}(M_{\lambda})\) for every \(\lambda \in S\) .

When A is a field, the simple A-modules are the vector spaces of dimension 1; every A-module is then semisimple (II, §7, No.~1, p.~292, Theorem 1), and its length is simply its dimension as a vector space over A (II, §7, No.~2, p.~293).

